% Options for packages loaded elsewhere
\PassOptionsToPackage{unicode}{hyperref}
\PassOptionsToPackage{hyphens}{url}
%
\documentclass[
]{article}
\usepackage{lmodern}
\usepackage{amssymb,amsmath}
\usepackage{ifxetex,ifluatex}
\ifnum 0\ifxetex 1\fi\ifluatex 1\fi=0 % if pdftex
  \usepackage[T1]{fontenc}
  \usepackage[utf8]{inputenc}
  \usepackage{textcomp} % provide euro and other symbols
\else % if luatex or xetex
  \usepackage{unicode-math}
  \defaultfontfeatures{Scale=MatchLowercase}
  \defaultfontfeatures[\rmfamily]{Ligatures=TeX,Scale=1}
\fi
% Use upquote if available, for straight quotes in verbatim environments
\IfFileExists{upquote.sty}{\usepackage{upquote}}{}
\IfFileExists{microtype.sty}{% use microtype if available
  \usepackage[]{microtype}
  \UseMicrotypeSet[protrusion]{basicmath} % disable protrusion for tt fonts
}{}
\makeatletter
\@ifundefined{KOMAClassName}{% if non-KOMA class
  \IfFileExists{parskip.sty}{%
    \usepackage{parskip}
  }{% else
    \setlength{\parindent}{0pt}
    \setlength{\parskip}{6pt plus 2pt minus 1pt}}
}{% if KOMA class
  \KOMAoptions{parskip=half}}
\makeatother
\usepackage{xcolor}
\IfFileExists{xurl.sty}{\usepackage{xurl}}{} % add URL line breaks if available
\IfFileExists{bookmark.sty}{\usepackage{bookmark}}{\usepackage{hyperref}}
\hypersetup{
  hidelinks,
  pdfcreator={LaTeX via pandoc}}
\urlstyle{same} % disable monospaced font for URLs
\usepackage[margin=1in]{geometry}
\usepackage{color}
\usepackage{fancyvrb}
\newcommand{\VerbBar}{|}
\newcommand{\VERB}{\Verb[commandchars=\\\{\}]}
\DefineVerbatimEnvironment{Highlighting}{Verbatim}{commandchars=\\\{\}}
% Add ',fontsize=\small' for more characters per line
\usepackage{framed}
\definecolor{shadecolor}{RGB}{248,248,248}
\newenvironment{Shaded}{\begin{snugshade}}{\end{snugshade}}
\newcommand{\AlertTok}[1]{\textcolor[rgb]{0.94,0.16,0.16}{#1}}
\newcommand{\AnnotationTok}[1]{\textcolor[rgb]{0.56,0.35,0.01}{\textbf{\textit{#1}}}}
\newcommand{\AttributeTok}[1]{\textcolor[rgb]{0.77,0.63,0.00}{#1}}
\newcommand{\BaseNTok}[1]{\textcolor[rgb]{0.00,0.00,0.81}{#1}}
\newcommand{\BuiltInTok}[1]{#1}
\newcommand{\CharTok}[1]{\textcolor[rgb]{0.31,0.60,0.02}{#1}}
\newcommand{\CommentTok}[1]{\textcolor[rgb]{0.56,0.35,0.01}{\textit{#1}}}
\newcommand{\CommentVarTok}[1]{\textcolor[rgb]{0.56,0.35,0.01}{\textbf{\textit{#1}}}}
\newcommand{\ConstantTok}[1]{\textcolor[rgb]{0.00,0.00,0.00}{#1}}
\newcommand{\ControlFlowTok}[1]{\textcolor[rgb]{0.13,0.29,0.53}{\textbf{#1}}}
\newcommand{\DataTypeTok}[1]{\textcolor[rgb]{0.13,0.29,0.53}{#1}}
\newcommand{\DecValTok}[1]{\textcolor[rgb]{0.00,0.00,0.81}{#1}}
\newcommand{\DocumentationTok}[1]{\textcolor[rgb]{0.56,0.35,0.01}{\textbf{\textit{#1}}}}
\newcommand{\ErrorTok}[1]{\textcolor[rgb]{0.64,0.00,0.00}{\textbf{#1}}}
\newcommand{\ExtensionTok}[1]{#1}
\newcommand{\FloatTok}[1]{\textcolor[rgb]{0.00,0.00,0.81}{#1}}
\newcommand{\FunctionTok}[1]{\textcolor[rgb]{0.00,0.00,0.00}{#1}}
\newcommand{\ImportTok}[1]{#1}
\newcommand{\InformationTok}[1]{\textcolor[rgb]{0.56,0.35,0.01}{\textbf{\textit{#1}}}}
\newcommand{\KeywordTok}[1]{\textcolor[rgb]{0.13,0.29,0.53}{\textbf{#1}}}
\newcommand{\NormalTok}[1]{#1}
\newcommand{\OperatorTok}[1]{\textcolor[rgb]{0.81,0.36,0.00}{\textbf{#1}}}
\newcommand{\OtherTok}[1]{\textcolor[rgb]{0.56,0.35,0.01}{#1}}
\newcommand{\PreprocessorTok}[1]{\textcolor[rgb]{0.56,0.35,0.01}{\textit{#1}}}
\newcommand{\RegionMarkerTok}[1]{#1}
\newcommand{\SpecialCharTok}[1]{\textcolor[rgb]{0.00,0.00,0.00}{#1}}
\newcommand{\SpecialStringTok}[1]{\textcolor[rgb]{0.31,0.60,0.02}{#1}}
\newcommand{\StringTok}[1]{\textcolor[rgb]{0.31,0.60,0.02}{#1}}
\newcommand{\VariableTok}[1]{\textcolor[rgb]{0.00,0.00,0.00}{#1}}
\newcommand{\VerbatimStringTok}[1]{\textcolor[rgb]{0.31,0.60,0.02}{#1}}
\newcommand{\WarningTok}[1]{\textcolor[rgb]{0.56,0.35,0.01}{\textbf{\textit{#1}}}}
\usepackage{longtable,booktabs}
% Correct order of tables after \paragraph or \subparagraph
\usepackage{etoolbox}
\makeatletter
\patchcmd\longtable{\par}{\if@noskipsec\mbox{}\fi\par}{}{}
\makeatother
% Allow footnotes in longtable head/foot
\IfFileExists{footnotehyper.sty}{\usepackage{footnotehyper}}{\usepackage{footnote}}
\makesavenoteenv{longtable}
\usepackage{graphicx}
\makeatletter
\def\maxwidth{\ifdim\Gin@nat@width>\linewidth\linewidth\else\Gin@nat@width\fi}
\def\maxheight{\ifdim\Gin@nat@height>\textheight\textheight\else\Gin@nat@height\fi}
\makeatother
% Scale images if necessary, so that they will not overflow the page
% margins by default, and it is still possible to overwrite the defaults
% using explicit options in \includegraphics[width, height, ...]{}
\setkeys{Gin}{width=\maxwidth,height=\maxheight,keepaspectratio}
% Set default figure placement to htbp
\makeatletter
\def\fps@figure{htbp}
\makeatother
\setlength{\emergencystretch}{3em} % prevent overfull lines
\providecommand{\tightlist}{%
  \setlength{\itemsep}{0pt}\setlength{\parskip}{0pt}}
\setcounter{secnumdepth}{5}
\usepackage{booktabs}
\hypersetup{colorlinks=true}

\author{}
\date{\vspace{-2.5em}}

\begin{document}

{
\setcounter{tocdepth}{2}
\tableofcontents
}
\hypertarget{r-basics}{%
\section{R Basics}\label{r-basics}}

\hypertarget{lets-get-started-with-rstudio}{%
\subsection{Let's get started with RStudio}\label{lets-get-started-with-rstudio}}

\begin{itemize}
\tightlist
\item
  Start installing \texttt{tidyverse}
\item
  Quick tour of RStudio's UI
\item
  Let's plot a thing (below)
\item
  Some idiosyncracies of my code

  \begin{itemize}
  \tightlist
  \item
    \texttt{\textquotesingle{}} rather than \texttt{"}
  \item
    \texttt{\#\#} for documentation comment
  \item
    \texttt{=} not \texttt{\textless{}-} for assignment
  \end{itemize}
\item
  Essential Global Options

  \begin{itemize}
  \tightlist
  \item
    Never save .Rdata, load workspace
  \end{itemize}
\end{itemize}

\begin{figure}

{\centering \includegraphics[width=0.6\linewidth]{images/Rstudio_config} 

}

\caption{Recommended RStudio configuration}\label{fig:unnamed-chunk-2}
\end{figure}

\begin{itemize}
\tightlist
\item
  Super useful keyboard commands:

  \begin{itemize}
  \tightlist
  \item
    Insert a pipe: ⌘+Shift+m
  \item
    Comment/uncomment: ⌘+Shift+c
  \item
    Tools \(\to\) Modify Keyboard Shortcuts\ldots{}

    \begin{itemize}
    \tightlist
    \item
      ⌘+1, ⌘+2, ⌘+3
    \end{itemize}
  \item
    Cheatsheet: \url{https://raw.githubusercontent.com/rstudio/cheatsheets/master/rstudio-ide.pdf}
  \end{itemize}
\item
  \textbf{If you're already familiar with the basics of R}, I recommend spending some time with @WickhamAdvanced2014.
\end{itemize}

\hypertarget{lets-plot-a-thing}{%
\subsubsection{Let's plot a thing}\label{lets-plot-a-thing}}

\begin{Shaded}
\begin{Highlighting}[]
\KeywordTok{library}\NormalTok{(tidyverse)}
\KeywordTok{library}\NormalTok{(gapminder)}
\KeywordTok{library}\NormalTok{(plotly)}

\KeywordTok{data}\NormalTok{(gapminder) }\CommentTok{\#\# Gets a dataset from a loaded package}
\NormalTok{gapminder}
\NormalTok{?gapminder}

\KeywordTok{ggplot}\NormalTok{(}\DataTypeTok{data =}\NormalTok{ gapminder, }\KeywordTok{aes}\NormalTok{(}\DataTypeTok{x =}\NormalTok{ gdpPercap, }\DataTypeTok{y =}\NormalTok{ lifeExp, }\DataTypeTok{color =}\NormalTok{ continent, }
                             \DataTypeTok{size =} \KeywordTok{log10}\NormalTok{(pop))) }\OperatorTok{+}
\StringTok{  }\KeywordTok{geom\_point}\NormalTok{() }\OperatorTok{+}
\StringTok{  }\KeywordTok{facet\_wrap}\NormalTok{(}\KeywordTok{vars}\NormalTok{(continent))}

\KeywordTok{ggplot}\NormalTok{(}\DataTypeTok{data =}\NormalTok{ gapminder, }\KeywordTok{aes}\NormalTok{(}\DataTypeTok{x =}\NormalTok{ year, }\DataTypeTok{y =}\NormalTok{ lifeExp, }\DataTypeTok{color =}\NormalTok{ continent)) }\OperatorTok{+}
\StringTok{  }\KeywordTok{geom\_line}\NormalTok{(}\KeywordTok{aes}\NormalTok{(}\DataTypeTok{group =}\NormalTok{ country)) }\OperatorTok{+}
\StringTok{  }\KeywordTok{facet\_wrap}\NormalTok{(}\KeywordTok{vars}\NormalTok{(continent))}

\KeywordTok{ggplotly}\NormalTok{()}
\end{Highlighting}
\end{Shaded}

\hypertarget{exercise}{%
\subsubsection{Exercise}\label{exercise}}

\begin{itemize}
\tightlist
\item
  Moar plots
\end{itemize}

\hypertarget{selecting-elements-and-some-math}{%
\subsection{Selecting elements, and some math}\label{selecting-elements-and-some-math}}

\begin{Shaded}
\begin{Highlighting}[]
\NormalTok{gapminder}\OperatorTok{$}\NormalTok{lifeExp}
\CommentTok{\#\# These are useful in certain circumstances, but harder to read}
\NormalTok{gapminder[[}\StringTok{\textquotesingle{}lifeExp\textquotesingle{}}\NormalTok{]]}
\NormalTok{gapminder[,}\StringTok{\textquotesingle{}lifeExp\textquotesingle{}}\NormalTok{]}
\NormalTok{gapminder[,}\DecValTok{4}\NormalTok{]}

\KeywordTok{mean}\NormalTok{(gapminder}\OperatorTok{$}\NormalTok{lifeExp)}
\KeywordTok{sd}\NormalTok{(gapminder}\OperatorTok{$}\NormalTok{lifeExp)}
\KeywordTok{log}\NormalTok{(gapminder}\OperatorTok{$}\NormalTok{lifeExp)}
\NormalTok{gapminder}\OperatorTok{$}\NormalTok{lifeExp }\OperatorTok{/}\StringTok{ }\DecValTok{100}
\NormalTok{gapminder}\OperatorTok{$}\NormalTok{year}\OperatorTok{\^{}}\DecValTok{2}

\KeywordTok{cor}\NormalTok{(gapminder}\OperatorTok{$}\NormalTok{lifeExp, gapminder}\OperatorTok{$}\NormalTok{year)}
\end{Highlighting}
\end{Shaded}

\hypertarget{exercise-1}{%
\subsubsection{Exercise}\label{exercise-1}}

Calculate correlations between each pair of life expectancy, population, and GDP per capita

\hypertarget{base-r-data-types}{%
\subsection{Base R data types}\label{base-r-data-types}}

\hypertarget{numerics}{%
\subsubsection{Numerics}\label{numerics}}

\begin{Shaded}
\begin{Highlighting}[]
\NormalTok{foo =}\StringTok{ }\DecValTok{2}
\NormalTok{bar =}\StringTok{ }\DecValTok{3}
\NormalTok{foobar =}\StringTok{ }\NormalTok{foo }\OperatorTok{+}\StringTok{ }\NormalTok{bar}
\NormalTok{foo }\OperatorTok{/}\StringTok{ }\NormalTok{bar}
\KeywordTok{try}\NormalTok{(foo }\OperatorTok{/}\StringTok{ }\DecValTok{0}\NormalTok{)}
\end{Highlighting}
\end{Shaded}

\hypertarget{vectors}{%
\subsubsection{Vectors}\label{vectors}}

\begin{Shaded}
\begin{Highlighting}[]
\NormalTok{foo =}\StringTok{ }\KeywordTok{c}\NormalTok{(}\DecValTok{2}\NormalTok{, }\DecValTok{3}\NormalTok{, }\DecValTok{4}\NormalTok{)}
\NormalTok{bar =}\StringTok{ }\KeywordTok{c}\NormalTok{(}\DecValTok{5}\NormalTok{, }\DecValTok{6}\NormalTok{, }\DecValTok{7}\NormalTok{)}
\NormalTok{foo[}\DecValTok{2}\NormalTok{]}
\NormalTok{foo[}\DecValTok{2}\OperatorTok{:}\DecValTok{3}\NormalTok{]}
\KeywordTok{try}\NormalTok{(foo[}\DecValTok{1}\NormalTok{, }\DecValTok{3}\NormalTok{])}
\NormalTok{foo[}\KeywordTok{c}\NormalTok{(}\DecValTok{1}\NormalTok{, }\DecValTok{3}\NormalTok{)]}

\NormalTok{foo }\OperatorTok{+}\StringTok{ }\NormalTok{bar}
\NormalTok{foo }\OperatorTok{*}\StringTok{ }\NormalTok{bar}
\CommentTok{\#\# NB difference between vector and length{-}1 comparisons}
\NormalTok{bar }\OperatorTok{==}\StringTok{ }\NormalTok{foo }\OperatorTok{+}\StringTok{ }\DecValTok{3}
\KeywordTok{identical}\NormalTok{(bar, foo }\OperatorTok{+}\StringTok{ }\DecValTok{3}\NormalTok{)}

\KeywordTok{c}\NormalTok{(foo, bar)}
\end{Highlighting}
\end{Shaded}

\hypertarget{strings}{%
\subsubsection{Strings}\label{strings}}

\begin{Shaded}
\begin{Highlighting}[]
\NormalTok{foo =}\StringTok{ }\KeywordTok{c}\NormalTok{(}\StringTok{\textquotesingle{}monkey\textquotesingle{}}\NormalTok{, }\StringTok{\textquotesingle{}banana\textquotesingle{}}\NormalTok{)}
\NormalTok{bar =}\StringTok{ }\KeywordTok{c}\NormalTok{(}\StringTok{\textquotesingle{}house\textquotesingle{}}\NormalTok{)}
\CommentTok{\#\# How to put these together? }
\KeywordTok{try}\NormalTok{(foo }\OperatorTok{+}\StringTok{ }\NormalTok{bar)}
\KeywordTok{c}\NormalTok{(foo, bar)}
\KeywordTok{paste}\NormalTok{(foo, bar)}
\KeywordTok{paste0}\NormalTok{(foo, bar)}
\end{Highlighting}
\end{Shaded}

\hypertarget{factors}{%
\subsubsection{Factors}\label{factors}}

\begin{Shaded}
\begin{Highlighting}[]
\NormalTok{foo\_str =}\StringTok{ }\KeywordTok{c}\NormalTok{(}\StringTok{\textquotesingle{}foo\textquotesingle{}}\NormalTok{, }\StringTok{\textquotesingle{}foo\textquotesingle{}}\NormalTok{, }\StringTok{\textquotesingle{}bar\textquotesingle{}}\NormalTok{, }\StringTok{\textquotesingle{}zoo\textquotesingle{}}\NormalTok{)}
\KeywordTok{as.integer}\NormalTok{(foo\_str)}
\NormalTok{foo\_fct =}\StringTok{ }\KeywordTok{as.factor}\NormalTok{(foo\_str)}
\KeywordTok{as.integer}\NormalTok{(foo)}
\KeywordTok{levels}\NormalTok{(foo\_fct)}
\end{Highlighting}
\end{Shaded}

\hypertarget{lists}{%
\subsubsection{Lists}\label{lists}}

\begin{Shaded}
\begin{Highlighting}[]
\NormalTok{some\_list =}\StringTok{ }\KeywordTok{list}\NormalTok{(}\StringTok{\textquotesingle{}foo\textquotesingle{}}\NormalTok{, }\StringTok{\textquotesingle{}bar\textquotesingle{}}\NormalTok{, }\DecValTok{12}\NormalTok{)}
\NormalTok{some\_list}
\NormalTok{some\_list[}\DecValTok{2}\NormalTok{]}

\NormalTok{nested\_list =}\StringTok{ }\KeywordTok{list}\NormalTok{(}\StringTok{\textquotesingle{}foo\textquotesingle{}}\NormalTok{, }\StringTok{\textquotesingle{}bar\textquotesingle{}}\NormalTok{, some\_list)}
\NormalTok{nested\_list}
\NormalTok{nested\_list[[}\DecValTok{3}\NormalTok{]]}

\NormalTok{moar\_list =}\StringTok{ }\KeywordTok{list}\NormalTok{(}\StringTok{\textquotesingle{}foo\textquotesingle{}}\NormalTok{ =}\StringTok{ }\DecValTok{12}\NormalTok{, }\StringTok{\textquotesingle{}bar\textquotesingle{}}\NormalTok{ =}\StringTok{ }\DecValTok{13}\NormalTok{)}
\NormalTok{moar\_list}\OperatorTok{$}\NormalTok{foo}
\end{Highlighting}
\end{Shaded}

\hypertarget{data-frames}{%
\subsubsection{Data frames}\label{data-frames}}

\begin{itemize}
\tightlist
\item
  \texttt{gapminder} is a data frame
\item
  Technical point: a data frame is a list

  \begin{itemize}
  \tightlist
  \item
    ``Columns'' = list elements
  \item
    Restriction: all of the elements must have the same length
  \end{itemize}
\item
  A tibble (\texttt{tbl\_df}) is a data frame with nice printing (part of the tidyverse)

  \begin{itemize}
  \tightlist
  \item
    Plus some technical restrictions that make it nicer to work with
  \end{itemize}
\end{itemize}

Like 99\% of the time you create a data frame by loading data or manipulating another data frame. But it's good to know how to create them by hand:

\begin{Shaded}
\begin{Highlighting}[]
\NormalTok{foo =}\StringTok{ }\DecValTok{1}\OperatorTok{:}\DecValTok{5}
\NormalTok{bar =}\StringTok{ }\NormalTok{letters[}\DecValTok{1}\OperatorTok{:}\DecValTok{5}\NormalTok{]}

\KeywordTok{data.frame}\NormalTok{(foo, bar)}
\KeywordTok{tibble}\NormalTok{(foo, bar)}

\CommentTok{\#\# tribble() is super useful when you\textquotesingle{}re creating a dataframe entirely by hand}
\KeywordTok{tribble}\NormalTok{(}\OperatorTok{\textasciitilde{}}\NormalTok{foo, }\OperatorTok{\textasciitilde{}}\NormalTok{bar, }
        \DecValTok{1}\NormalTok{, }\StringTok{\textquotesingle{}a\textquotesingle{}}\NormalTok{, }
        \DecValTok{2}\NormalTok{, }\StringTok{\textquotesingle{}b\textquotesingle{}}\NormalTok{, }
        \DecValTok{3}\NormalTok{, }\StringTok{\textquotesingle{}c\textquotesingle{}}\NormalTok{, }
        \DecValTok{4}\NormalTok{, }\StringTok{\textquotesingle{}d\textquotesingle{}}\NormalTok{, }
        \DecValTok{5}\NormalTok{, }\StringTok{\textquotesingle{}e\textquotesingle{}}\NormalTok{)}
\end{Highlighting}
\end{Shaded}

\hypertarget{examining-variables}{%
\subsection{Examining variables}\label{examining-variables}}

\begin{itemize}
\tightlist
\item
  \texttt{length()}, \texttt{dim()}
\item
  \texttt{str()}
\item
  \texttt{class()}
\item
  \texttt{head()}
\item
  \href{https://gge-ucd.github.io/R-DAVIS/lesson_how_r_thinks_about_data.html\#challenge}{challenge}
\end{itemize}

\hypertarget{exercise-2}{%
\subsubsection{Exercise}\label{exercise-2}}

\begin{itemize}
\tightlist
\item
  What's the difference between \texttt{gapminder{[}\textquotesingle{}country\textquotesingle{}{]}} (single brackets) and \texttt{gapminder{[}{[}\textquotesingle{}country\textquotesingle{}{]}{]}} (double brackets)?
\end{itemize}

\hypertarget{logical-operators}{%
\subsection{Logical operators}\label{logical-operators}}

\begin{Shaded}
\begin{Highlighting}[]
\ControlFlowTok{if}\NormalTok{ (}\OtherTok{TRUE}\NormalTok{) \{}
  \KeywordTok{print}\NormalTok{(}\StringTok{\textquotesingle{}it}\CharTok{\textbackslash{}\textquotesingle{}}\StringTok{s true!\textquotesingle{}}\NormalTok{)}
\NormalTok{\} }\ControlFlowTok{else}\NormalTok{ \{}
  \KeywordTok{print}\NormalTok{(}\StringTok{\textquotesingle{}it}\CharTok{\textbackslash{}\textquotesingle{}}\StringTok{s false!\textquotesingle{}}\NormalTok{)}
\NormalTok{\}}

\NormalTok{foo =}\StringTok{ }\KeywordTok{c}\NormalTok{(}\DecValTok{1}\NormalTok{, }\DecValTok{{-}1}\NormalTok{, }\DecValTok{{-}3}\NormalTok{)}
\NormalTok{foo }\OperatorTok{\textgreater{}}\StringTok{ }\DecValTok{0}
\ControlFlowTok{if}\NormalTok{ (foo }\OperatorTok{\textgreater{}}\StringTok{ }\DecValTok{0}\NormalTok{) \{}
  \KeywordTok{print}\NormalTok{(}\StringTok{\textquotesingle{}greater than zero!\textquotesingle{}}\NormalTok{)}
\NormalTok{\}}

\ControlFlowTok{if}\NormalTok{ (}\KeywordTok{any}\NormalTok{(foo }\OperatorTok{\textgreater{}}\StringTok{ }\DecValTok{0}\NormalTok{)) \{}
  \KeywordTok{print}\NormalTok{(}\StringTok{\textquotesingle{}at least one is greater than zero!\textquotesingle{}}\NormalTok{)}
\NormalTok{\}}
\ControlFlowTok{if}\NormalTok{ (}\KeywordTok{all}\NormalTok{(foo }\OperatorTok{\textgreater{}}\StringTok{ }\DecValTok{0}\NormalTok{)) \{}
  \KeywordTok{print}\NormalTok{(}\StringTok{\textquotesingle{}all are greater than zero!\textquotesingle{}}\NormalTok{)}
\NormalTok{\} }\ControlFlowTok{else}\NormalTok{ \{}
  \KeywordTok{print}\NormalTok{(}\StringTok{\textquotesingle{}at least one is not greater than zero!\textquotesingle{}}\NormalTok{)}
\NormalTok{\}}

\ControlFlowTok{if}\NormalTok{ (foo }\OperatorTok{\textgreater{}}\StringTok{ }\DecValTok{0} \OperatorTok{|}\StringTok{ }\NormalTok{foo }\OperatorTok{\textless{}}\StringTok{ }\DecValTok{0}\NormalTok{) \{}
  \KeywordTok{print}\NormalTok{(}\StringTok{\textquotesingle{}non{-}zero\textquotesingle{}}\NormalTok{)}
\NormalTok{\}}
\ControlFlowTok{if}\NormalTok{ (foo }\OperatorTok{\textgreater{}}\StringTok{ }\DecValTok{0} \OperatorTok{||}\StringTok{ }\NormalTok{foo }\OperatorTok{\textless{}}\StringTok{ }\DecValTok{0}\NormalTok{) \{}
  \KeywordTok{print}\NormalTok{(}\StringTok{\textquotesingle{}non{-}zero\textquotesingle{}}\NormalTok{)}
\NormalTok{\}}

\ControlFlowTok{if}\NormalTok{ (foo }\OperatorTok{!=}\StringTok{ }\DecValTok{0}\NormalTok{) \{}
  \KeywordTok{print}\NormalTok{(}\StringTok{\textquotesingle{}non{-}zero\textquotesingle{}}\NormalTok{)}
\NormalTok{\}}
\ControlFlowTok{if}\NormalTok{ (}\KeywordTok{all}\NormalTok{(foo }\OperatorTok{!=}\StringTok{ }\DecValTok{0}\NormalTok{)) \{}
  \KeywordTok{print}\NormalTok{(}\StringTok{\textquotesingle{}non{-}zero\textquotesingle{}}\NormalTok{)}
\NormalTok{\}}

\NormalTok{bar =}\StringTok{ }\KeywordTok{c}\NormalTok{(}\DecValTok{17}\NormalTok{, }\DecValTok{25}\NormalTok{, }\DecValTok{0}\NormalTok{, }\DecValTok{{-}14}\NormalTok{)}
\KeywordTok{ifelse}\NormalTok{(bar }\OperatorTok{\textgreater{}}\StringTok{ }\DecValTok{0}\NormalTok{, }\StringTok{\textquotesingle{}positive\textquotesingle{}}\NormalTok{, }\StringTok{\textquotesingle{}negative or zero\textquotesingle{}}\NormalTok{)}
\KeywordTok{case\_when}\NormalTok{(bar }\OperatorTok{\textgreater{}}\StringTok{ }\DecValTok{0}  \OperatorTok{\textasciitilde{}}\StringTok{ \textquotesingle{}positive\textquotesingle{}}\NormalTok{, }
\NormalTok{          bar }\OperatorTok{==}\StringTok{ }\DecValTok{0} \OperatorTok{\textasciitilde{}}\StringTok{ \textquotesingle{}zero\textquotesingle{}}\NormalTok{, }
\NormalTok{          bar }\OperatorTok{\textless{}}\StringTok{ }\DecValTok{0}  \OperatorTok{\textasciitilde{}}\StringTok{ \textquotesingle{}negative\textquotesingle{}}\NormalTok{)}
\end{Highlighting}
\end{Shaded}

\hypertarget{missing-data}{%
\subsection{Missing data}\label{missing-data}}

\url{https://gge-ucd.github.io/R-DAVIS/lesson_how_r_thinks_about_data.html\#missing_data}

\begin{Shaded}
\begin{Highlighting}[]
\NormalTok{missing =}\StringTok{ }\KeywordTok{c}\NormalTok{(}\DecValTok{7}\NormalTok{, }\DecValTok{17}\NormalTok{, }\OtherTok{NA}\NormalTok{, }\DecValTok{12}\NormalTok{)}
\NormalTok{missing}
\KeywordTok{mean}\NormalTok{(missing)}
\KeywordTok{mean}\NormalTok{(missing, }\DataTypeTok{na.rm =} \OtherTok{TRUE}\NormalTok{)}
\KeywordTok{is.na}\NormalTok{(missing)}
\end{Highlighting}
\end{Shaded}

\hypertarget{exercise-3}{%
\subsubsection{Exercise}\label{exercise-3}}

\begin{itemize}
\tightlist
\item
  What do the logical operators do with \texttt{NA}s?
\end{itemize}

\end{document}
